\documentclass[11pt]{article}
\usepackage[a4paper,margin=26mm]{geometry}
\usepackage{amsmath,amssymb}
\usepackage[hidelinks]{hyperref}
\setlength{\emergencystretch}{3em}
\title{Conjecture 00000000684: a classical circle has an infinite finitely generated group}
\author{lizaixi01, with OpenAI Codex assistance}
\date{23 October 2026}
\begin{document}
\maketitle
\section{The standard virtual group and the counterexample}
We use the standard Kauffman group $G(D)$ of a virtual diagram,
defined by its Wirtinger presentation. The source names a virtual
diagram group without further specifying a variant. This standard
model is defined in Kauffman,
\href{https://homepages.math.uic.edu/~kauffman/VKT.pdf}
{\emph{Virtual Knot Theory}}, Europ. J. Combinatorics 20 (1999),
Section 4, printed pp.671--672: generators correspond to arcs between
classical undercrossings, and relations to classical crossings.
Virtual crossings contribute no new generators or relations.
At printed p.663 the author explicitly uses ``trivial fundamental
group'' to mean the infinite cyclic group $\mathbf Z$.

Take the actual round circle
\[
 S^1=\{z\in\mathbf C:|z|=1\},\qquad
 K:S^1\longrightarrow\mathbf C\times\mathbf R=\mathbf R^3,
 \quad K(z)=(z,0).
\]
Its planar projection is $z\mapsto z$. It is an actual embedding,
with no self-crossings and no virtual crossings. Thus it gives a
classical diagram of the ordinary round unknot. Its standard
Wirtinger group has one arc generator and no crossing relations,
so is isomorphic to $\mathbf Z$. In particular it is finitely
generated and infinite, although this knot is classical. This
contradicts the forward implication of the proposed equivalence.
An additional Gauss-diagram criterion cannot repair that false
necessary implication.

\section{Actual crossings and arcs, rather than an assigned presentation}
The formal source uses Mathlib's actual complex unit circle.
It first considers a crossing-free diagram given by a continuous
injective map $f:S^1\to\mathbf C$. The set of crossing preimages is
defined from the actual map:
\[
 U_f=\{x\in S^1:\exists y\ne x,\ f(y)=f(x)\}.
\]
Injectivity proves $U_f=\varnothing$. Every possible undercrossing
preimage subset is therefore empty as well. This avoids treating
overcrossing preimages as undercrossings in a general diagram;
the two sets are both empty in this subclass.

Arcs are the genuine connected components of $S^1\setminus U_f$.
An explicit homeomorphism identifies this actual complement with
$S^1$, and the connectedness of the circle proves that its component
type has exactly one element. The formal generator type is this
component type, not a separately assigned one-element label.

\newpage
The classical-crossing index type is also defined from actual
pairs of distinct circle points with equal planar image. It is
proved empty. A crossing-relator assignment has this type as its
domain, and its range is consequently empty. Lean proves this for
\emph{every} relation assignment on the actual empty crossing type.
Thus no choice of signs, underpassing branch or usual Wirtinger
convention can add a crossing relation in this example. There is
no need to encode the crossing formula for nonempty diagrams to
settle this crossing-free case.

The diagram group is the actual \texttt{PresentedGroup} on those
arc generators and the crossing-indexed relator set. Lean proves
that its quotient map from the free group is injective as well as
surjective, yielding the genuine isomorphisms
\[
 G(D)\simeq F(\text{one actual arc})\simeq (\mathbf Z,+).
\]
Mathlib's actual group finite-generation property and infinitude
are then proved for this presented group.

\section{Classical realization and the exact scope}
For the round diagram, the projection is the actual subtype inclusion
$S^1\hookrightarrow\mathbf C$. The formal source also constructs
the actual spatial embedding $z\mapsto(f(z),0)$ and checks that its
projection is $f$ and its height is zero. For $f(z)=z$ it additionally
records the unit-norm equation, binding the construction to the
ordinary round circle.

The predicate \texttt{ClassicalRealization} records this sufficient
classical realization for the crossing-free subclass. It is not
presented as a classification or definition of arbitrary virtual
knots. In particular, it does not replace nonclassicality for
arbitrary diagrams by failure to embed in a plane. The standard
round diagram already has no virtual crossings, so its classical
status requires no equivalence or unknot-recognition algorithm.

The final formal statements refute both the restricted equivalence
and, separately, its necessary implication: an infinite finitely
generated diagram group would prohibit such an explicit classical
realization. The actual round diagram satisfies the group properties
and has that realization, so this implication fails.

We use Kauffman's named standard Wirtinger model. The Lean source
does not reprove invariance under every extended Reidemeister move,
the identification with the fundamental group of a complement,
or a theorem about every alternative extended virtual-group
definition. None of these additional results is used to assign a
group to an unrelated presentation: the arc and crossing types are
derived from the actual circle projection before forming the group.

\section{Verification}
Lean 4.33.0 and Mathlib commit
\begin{center}\small\texttt{db584cd6d46c92f209a44c0f1c829460d327499d}\end{center}
are pinned. The package includes a new root build, transitive axiom
audits and full official \texttt{leanchecker --fresh Main} replay.
No \texttt{sorry}, \texttt{native\_decide} or extra axiom is used.
OpenAI Codex assisted; a separate supervising Codex agent reviewed
source, report and PDF. This is AI review, not human certification.
Organizer acceptance is pending.
\end{document}
